\chapter{Prompting Multimodal: Texto, Imagem, Áudio e Ferramentas}
\label{ch:multimodal}

\section{O que quer dizer ``multimodal'' na prática}

“Multimodal” é uma daquelas palavras que toda a gente usa e poucos definem. De forma pragmática:

\begin{itemize}[leftmargin=*]
\item \textbf{Modelo unimodal}:
\begin{itemize}[leftmargin=*]
\item trabalha com uma modalidade, que pode ser texto, imagem, áudio ou outra,
\item outras modalidades exigem componentes ou conversões adicionais.
\end{itemize}
\item \textbf{Modelo multimodal}:
\begin{itemize}[leftmargin=*]
\item aceita mais do que um tipo de input (texto + imagem, texto + áudio, etc.),
\item pode aceitar ou produzir diferentes modalidades; aceitar imagem não implica gerar imagem,
\item pode combinar isto com ferramentas externas.
\end{itemize}
\end{itemize}

Em vez de pensares “modelo mágico que vê e ouve”, pensa:

\begin{itemize}[leftmargin=*]
\item podes ter modelos com representações partilhadas entre modalidades ou cadeias de componentes especializados. As capacidades exatas dependem da arquitetura~\citep{gpt4report2023}. A aplicação pode:
\begin{itemize}[leftmargin=*]
\item extrair informação de uma imagem (como se alguém a descrevesse em detalhe),
\item gerar uma imagem a partir de uma descrição,
\item transformar áudio em texto (ASR) e texto em áudio (TTS),
\item coordenar isto com outras APIs.
\end{itemize}
\end{itemize}

O teu trabalho, como engenheiro de prompts, é:

\begin{itemize}[leftmargin=*]
\item formular pedidos que usem bem esses canais (sem superegoísmo: “faz tudo sozinho”),
\item combinar etapas de texto+imagem+áudio em fluxos coerentes,
\item respeitar limites de cada modalidade (resolução, ruído, ambiguidade).
\end{itemize}

\begin{keyideabox}{Multimodal $\neq$ licença para preguiça}
Meter um screenshot e dizer “resolve isto” é pouco melhor do que mandar uma foto do quadro ao professor e exigir nota 20.
Continua a ser preciso dizer \emph{o que queres} que o modelo faça com aquilo que está na imagem/áudio.
\end{keyideabox}

\section{Usar imagens como input: screenshots, gráficos e documentos}

Modelos com visão conseguem “ler”:

\begin{itemize}[leftmargin=*]
\item screenshots de aplicações,
\item gráficos e dashboards,
\item fotografias de quadros brancos,
\item páginas de documentos (PDFs, apresentações).
\end{itemize}

Mas a forma como pedes faz muita diferença.

\subsection*{Padrões para screenshots de aplicações}

Mau prompt:

\begin{lstlisting}
[envias screenshot de um erro]
O que se passa aqui?
\end{lstlisting}

Melhor:

\begin{lstlisting}
[envias screenshot]

Contexto:

Aplicação: painel de gestão interno.

Ação que fiz: cliquei em "Guardar" após editar
os campos [A, B, C].

Expectativa: os dados deveriam ser atualizados
sem erro.

Pedido:

Descreve o que vês na imagem relevante para o erro.

Identifica mensagens de erro ou sinais de problema.

Propõe 3 hipóteses possíveis para a causa do erro
(lado cliente, servidor, dados).

Sugere próximos passos de debug.
\end{lstlisting}

Repara que:

\begin{itemize}[leftmargin=*]
\item dás contexto que a imagem não tem,
\item divides o pedido em sub-tarefas,
\item focas o modelo em “diagnosticar” em vez de “adivinhar”.
\end{itemize}

\subsection*{Padrões para gráficos e dashboards}

Quando é um gráfico de dados:

\begin{lstlisting}
[envias imagem de um gráfico ou dashboard]

Quero que assumes o papel de analista de dados.

Para a imagem:

Descreve sucintamente o que o gráfico mostra
(variáveis, eixos, período).

Identifica tendências principais (subida, descida,
sazonalidade, outliers).

Faz 3 hipóteses possíveis para explicar a tendência
principal.

Sugere 2-3 análises adicionais que deveríamos fazer
com os dados originais (não disponíveis na imagem).
\end{lstlisting}

Se não disseres que é um gráfico, o modelo normalmente percebe sozinho.
Indicar a tarefa permite avaliar se a descrição, interpretação ou crítica cumpre o objetivo.

\subsection*{Padrões para documentos e slides}

Para páginas de relatório, slides, etc.:

\begin{lstlisting}
[envias uma ou mais imagens/páginas]

Quero um resumo focado em:

objetivo deste documento (1-2 frases);

principais conclusões (bullet points);

qualquer recomendação explícita que apareça;

campos ou números que possam ter impacto
financeiro/legal.

Depois, indica-me:

3 perguntas que um gestor deveria fazer após ler isto.
\end{lstlisting}

\begin{promptlabbox}{Modo ``OCR crítico''}
\textbf{Objetivo e contexto:} Transcrever uma página autorizada com texto ou fórmulas e conservar a imagem original.

Experimenta usar a visão do modelo como OCR+analista:

\begin{itemize}[leftmargin=*]
\item tira uma fotografia de um quadro com fórmulas ou de um slide confuso,
\item pede ao modelo para:
\begin{itemize}[leftmargin=*]
\item transcrever as equações/texto,
\item apontar ambiguidade ou problemas de formatação,
\item sugerir uma versão mais clara (por exemplo, em LaTeX).
\end{itemize}
\end{itemize}

Isto transforma “foto de apontamentos” em algo pesquisável e reutilizável.

\textbf{Resultado esperado e validação:} Confere caracteres, sinais, unidades e estrutura contra a região de origem. Marca trechos ilegíveis em vez de os completar por adivinhação.

\textbf{Falhas possíveis:} Aceitar uma equação plausível mas diferente ou tratar o texto extraído como fonte infalível.
\end{promptlabbox}

\section{Gerar imagens com prompts de texto}

Quando o modelo também gera imagens, voltamos à essência: \emph{especificar bem}. Pedidos vagos produzem:

\begin{itemize}[leftmargin=*]
\item imagens genéricas,
\item estilos cliché,
\item detalhes incoerentes.
\end{itemize}

Elementos a controlar num prompt de imagem:

\begin{itemize}[leftmargin=*]
\item \textbf{Sujeito principal}: quem/quê está em foco?
\item \textbf{Ação/pose}: o que está a acontecer?
\item \textbf{Ambiente}: interior/exterior, contexto?
\item \textbf{Estilo}: realista, desenho, vector, “flat”, etc.
\item \textbf{Composição}: close-up, plano geral, perspectiva.
\item \textbf{Cores/clima}: claro/escuro, minimalista, saturado.
\end{itemize}

Exemplo fraco:

\begin{lstlisting}
Desenha um dashboard bonito sobre vendas.
\end{lstlisting}

Exemplo melhor:

\begin{lstlisting}
Gera uma imagem raster com aspeto de ilustracao plana de um dashboard
de vendas numa aplicação web moderna.

Requisitos:

Fundo claro, layout limpo.

Um grafico ilustrativo, sem dados reais nem numeros
que possam ser confundidos com metricas de negocio.

Uma tabela pequena à direita com 5 produtos top.

Cores suaves, sem elementos 3D.
\end{lstlisting}

\begin{warningboxenv}{Imagens com pessoas reais, logotipos e marcas}
Gerar ou editar imagens com:
\begin{itemize}[leftmargin=*]
\item rostos de pessoas reais,
\item logotipos registados,
\item marcas comerciais,
\end{itemize}
tem implicações legais e éticas:

\begin{itemize}[leftmargin=*]
\item direitos de imagem,
\item direitos de autor,
\item risco de deepfakes e desinformação.
\end{itemize}

Mesmo que a ferramenta permita tecnicamente, pensa se \emph{deves}.
E verifica políticas de uso da plataforma e da tua organização.
\end{warningboxenv}

\section{Combinar texto + imagem em fluxos de trabalho}

A parte interessante é quando combinas:

\begin{itemize}[leftmargin=*]
    \item texto $\Rightarrow$ imagem,
    \item imagem $\Rightarrow$ texto,
    \item texto + imagem $\Rightarrow$ outra coisa (código, plano, crítica).
\end{itemize}


Alguns padrões úteis:

\subsection*{Design assistido}

\begin{enumerate}[leftmargin=*]
\item Desenhas um wireframe à mão, tiras foto.
\item Pedes ao modelo para:
\begin{itemize}[leftmargin=*]
\item descrever o layout,
\item avaliar problemas de usabilidade,
\item sugerir melhorias.
\end{itemize}
\item Com base nisso, pedes uma imagem mais refinada ou código HTML/CSS.
\end{enumerate}

\subsection*{Debug visual}

\begin{enumerate}[leftmargin=*]
\item Screenshot de um erro numa UI.
\item Modelo analisa elementos visuais + texto.
\item Sugere hipóteses de bug (CSS, lógica, layout responsivo).
\end{enumerate}

\subsection*{Documentação ``visual-first''}

\begin{enumerate}[leftmargin=*]
\item Tiras screenshots de ecrãs importantes da tua app.
\item O modelo gera:
\begin{itemize}[leftmargin=*]
\item descrições,
\item bullet points de funcionalidade,
\item secções de ajuda “como fazer X”.
\end{itemize}
\end{enumerate}

\begin{promptlabbox}{Pipeline texto $\Rightarrow$  imagem $\Rightarrow$  texto}
\textbf{Objetivo e contexto:} Observar perdas ao converter descrição em maquete raster e voltar a extrair estrutura.

Faz um teste:

\begin{enumerate}[leftmargin=*]
\item Descreve um layout de página (texto).
\item Gera uma imagem de mockup com o modelo.
\item Depois, alimenta a imagem ao modelo e pede-lhe para:
\begin{itemize}[leftmargin=*]
\item extrair a estrutura (componentes, hierarquia),
\item propor melhorias de UX.
\end{itemize}
\end{enumerate}

Repara em quanta informação se perde/ganha em cada conversão.
Isto ajuda-te a calibrar expectativas sobre “fidelidade” multimodal.

\textbf{Resultado esperado e validação:} Define previamente os elementos obrigatórios. Verifica presença, hierarquia e detalhes; a imagem é uma maquete sem valor de gráfico de dados.

\textbf{Falhas possíveis:} Usar números inventados como métricas, presumir que raster é vetor ou avaliar apenas beleza.
\end{promptlabbox}

\section{Áudio: ASR, TTS e comandos por voz}

Uma arquitetura comum de voz encadeia três peças; outras processam áudio diretamente sem uma transcrição como etapa obrigatória:

\begin{itemize}[leftmargin=*]
\item \textbf{ASR (Automatic Speech Recognition)}: converter fala em texto.
\item \textbf{LLM}: processar o texto (como sempre).
\item \textbf{TTS (Text-to-Speech)}: converter a resposta textual em voz.
\end{itemize}

Padrões comuns:

\subsection*{Resumos de reuniões}

\begin{enumerate}[leftmargin=*]
\item Gravas uma reunião.
\item ASR gera transcrição (com erros, “hums”, etc.).
\item LLM:
\begin{itemize}[leftmargin=*]
\item limpa a transcrição (se quiseres),
\item extrai decisões, ações, riscos.
\end{itemize}
\end{enumerate}

Prompt típico:

\begin{lstlisting}
A partir da transcrição abaixo de uma reunião,
organiza em:

Decisões tomadas

Ações (com responsável se indicado)

Prazos mencionados

Questões em aberto

Riscos

Português de Portugal, formato em bullet points.
Ignora "hums", small talk e interrupções irrelevantes.
\end{lstlisting}

\subsection*{Comandos por voz}

Fluxo genérico:

\begin{enumerate}[leftmargin=*]
\item Utilizador fala.
\item ASR converte para texto.
\item LLM interpreta comando e decide ação:
\begin{itemize}[leftmargin=*]
\item responder com texto,
\item chamar ferramenta,
\item pedir clarificação.
\end{itemize}
\end{enumerate}

Convém ter prompts de “interpretação robusta” que:

\begin{itemize}[leftmargin=*]
\item tolerem erros de transcrição,
\item clarifiquem ambiguidade (“queres A ou B?”),
\item não façam ações perigosas com base em frases meio cortadas.
\end{itemize}

\begin{warningboxenv}{Privacidade em áudio é ainda mais sensível}
Gravar conversas, reuniões e comandos de voz implica:

\begin{itemize}[leftmargin=*]
\item verificação da licitude da gravação e informação aos participantes, considerando também as regras nacionais aplicáveis,
\item decisões sobre retenção (quanto tempo guardas?),
\item risco de capturar informações que ninguém queria que fossem processadas por IA.
\end{itemize}

Não é só uma questão técnica; é legal e ética.
\end{warningboxenv}

\section{Ferramentas multimodais e agentes ``sensoriais''}

Quando combinas:

\begin{itemize}[leftmargin=*]
\item visão,
\item texto,
\item áudio,
\item ferramentas (APIs),
\end{itemize}

começas a ter agentes que:

\begin{itemize}[leftmargin=*]
\item “olham” para o estado de um sistema (screenshots),
\item “ouvem” comandos,
\item consultam dados,
\item respondem por texto/voz.
\end{itemize}

O papel do prompt aqui é:

\begin{itemize}[leftmargin=*]
\item definir que tipo de informação o modelo deve tirar de cada modalidade;
\item especificar a ordem de operações (ver imagem \(\to\) consultar API \(\to\) explicar resultado);
\item impor limites (“não executes ações destrutivas sem autorização explícita”).
\end{itemize}

Exemplo conceptual:

\begin{lstlisting}
És um assistente que ajuda a diagnosticar problemas
num painel de controlo industrial.

Inputs:

Texto do operador (perguntas, descrições).

Screenshot do painel.

Tarefas:

Descrever o estado atual com base nos elementos
visíveis (indicadores, alarmes, gráficos).

Cruzar isso com as regras abaixo (manual de operação).

Explicar ao operador, em linguagem simples,
o que parece estar a acontecer e propondo
2-3 ações seguras de diagnóstico.

Nunca sugiras ações que envolvam desligar sistemas
críticos sem confirmação adicional explícita.
\end{lstlisting}

\section{Limites e riscos específicos do multimodal}

Com multimodal, ganhas poder, mas também:

\begin{itemize}[leftmargin=*]
\item \textbf{Ambiguidade visual}: imagens podem ser pouco claras, cortadas, desfocadas.
\item \textbf{Leitura enganadora}: o modelo pode interpretar um detalhe visual de forma errada e construir uma narrativa plausível mas falsa.
\item \textbf{Reconhecimento de pessoas}: risco de identificação, vigilância, perfis sensíveis.
\item \textbf{Dependência da qualidade do ASR}: erros de transcrição podem distorcer significados importantes.
\end{itemize}

Boas práticas:

\begin{itemize}[leftmargin=*]
\item sempre que o output depender \emph{fortemente} de um detalhe visual, pede ao modelo para indicar o que viu (“dizes que o botão está cinzento; que label tem?”);
\item evita decisões automatizadas em cima de reconhecimento facial ou de emoções;
\item para reuniões, considera partilhar resumos com os participantes para revisão humana (não assumir que “se a IA resumiu, está certo”);
\item documenta como as imagens/áudios são recolhidos, processados e descartados.
\end{itemize}

\begin{keyideabox}{Multimodal é uma ferramenta de contexto, não um oráculo}
Usa imagens, áudio e afins para dar \emph{mais contexto} ao modelo, não para substituir completamente o teu julgamento.
Se uma decisão é crítica, valida sempre com humanos e dados diretos.
\end{keyideabox}

\section{Separar observação, interpretação e ação}
Num teste de interface, uma captura permite observar que um botão está cortado. Não permite concluir, sozinha, que uma atualização à base de dados falhou. Regista «observado na imagem», «hipótese» e «verificação necessária». Faz o mesmo com áudio: uma palavra transcrita com erro pode inverter uma instrução.

Para gráficos que comuniquem dados, usa ferramentas que desenhem os valores originais em SVG, PDF ou outro formato verificável. Uma imagem gerada de um dashboard é uma maquete, não evidência estatística. Para documentos, preserva a página ou região de origem para confirmar números e negações.
